% Source: https://artofproblemsolving.com/community/c3h1362072_bisectors
% Stable identity: aops-topic-1362072-q01
Both lines pass through the origin, so the bisectors do too. A point $(x,y)$ lies on an angle bisector exactly when it is equidistant from the two lines $x-y=0$ and $7x-y=0$:
\[
\dfrac{|x-y|}{\sqrt{2}}=\dfrac{|7x-y|}{\sqrt{50}}\iff 5|x-y|=|7x-y| .
\]
If $5(x-y)=7x-y$, then $-2x=4y$, giving $y=-\dfrac{1}{2}x$. If $5(x-y)=-(7x-y)$, then $12x=6y$, giving $y=2x$.

\begin{asy}
usepackage("amsmath");
size(6cm);
defaultpen(fontsize(10pt));
import graph;
real L=1.25;
draw((-L,-L)--(L,L),linewidth(0.9)); draw((-L/7,-L)--(L/7,L),linewidth(0.9));
draw((-0.62,-1.24)--(0.62,1.24),red+dashed); draw((-1.4,0.7)--(1.4,-0.7),blue+dashed);
draw((-1.45,0)--(1.45,0),gray,Arrow(4)); draw((0,-1.3)--(0,1.35),gray,Arrow(4));
label("$x$",(1.45,0),S,gray); label("$y$",(0,1.35),W,gray);
label("$y=x$",(L,L),E); label("$y=7x$",(-L/7,-L),S);
label("$y=2x$",(0.62,1.24),N,red); label("$y=-\dfrac{1}{2}x$",(1.4,-0.7),E,blue);
dot((0,0));
label("\mbox{}",(min(currentpicture,user=true)+max(currentpicture,user=true))/2,black); // keeps the text after the figure black when LaTeX typesets the labels inline
\end{asy}

So the bisectors are $\boxed{y=2x\ \text{ and }\ y=-\dfrac{1}{2}x}$; they are perpendicular, as the two bisectors of a pair of lines must be. The line $y=2x$ bisects the acute angle between $y=x$ and $y=7x$. Check: the angle from $y=x$ to $y=2x$ has $\tan=\dfrac{2-1}{1+2}=\dfrac{1}{3}$, and from $y=2x$ to $y=7x$ has $\tan=\dfrac{7-2}{1+14}=\dfrac{1}{3}$.
